\centerline{\bf \Huge Многочлены над конечным полем}

\textbf{Определение.} Множество элементов $F$ с введёнными на нём алгебраическими операциями сложения + и умножения $*(\forall a, b \in F \quad(a+b) \in F, a * b \in F)$ называется полем ( $F,+, *$ ), если выполнены следующие аксиомы:

\begin{itemize}
    \item Коммутативность сложения: $\forall a, b \in F \quad a+b=b+a$.
    
    \item Ассоциативность сложения: $\forall a, b, c \in F \quad(a+b)+c=a+(b+c)$.
    
    \item Существование нулевого элемента: $\exists 0 \in F: \forall a \in F \quad a+0=a$.
    
    \item Существование противоположного элемента: $\forall a \in F \exists(-a) \in F: a+(-a)=0$.
    
    \item Коммутативность умножения: $\forall a, b \in F \quad a * b=b * a$.
    
    \item Ассоциативность умножения: $\forall a, b, c \in F \quad(a * b) * c=a *(b * c)$.
    
    \item Существование единичного элемента: $\exists 1 \in F: \forall a \in F \quad a * 1=a$.
    
    \item Существование обратного элемента для ненулевых элементов: $(\forall a \in F: a \neq 0) \exists a^{-1} \in F: a * a^{-1}=1$.
    
    \item Дистрибутивность умножения относительно сложения: $\forall a, b, c \in F \quad(a+b) * c=(a * c)+(b * c)$.
\end{itemize}


\problem{1}
Докажите, что множество остатков при делении на простое число $p$ является полем. Оно обозначается $\mathbb{Z}_p$.

\textbf{Определение.} Многочленом $f(x)$ над конечным полем $\mathbb{F}$ называется формальная сумма вида

$$
f(x)=f_0+f_1 x+\ldots+f_m x^m, f_i \in \mathbb{F}, f_m \neq 0 .
$$


Множество многочленов над полем $\mathbb{F}$ обозначается $\mathbb{F}[x]$.

\problem{2}
(a) Сформулируйте и докажите теорему Безу для многочленов над полем $\mathbb{Z}_p$.

(b) Сформулируйте и докажите теорему Виета для многочленов над полем $\mathbb{Z}_p$.

\problem{3} 
(a) Разложите на множители над $\mathbb{Z}_p$ многочлен $x^{p-1}-1$.

(b) Пользуясь предыдущим пунктом, докажите теорему Вильсона:
$(p-1)!\equiv-1(\bmod p)$ при простом $p$.

\problem{4}
\textit{(Критерий Эйзенштейна)}
Пусть $f(x)$ - многочлен с целыми, у которого старший коэффициент не делится на простое число $p$, все остальные коэффициенты делятся на $p$, а свободный член не делится на $p^2$. Тогда $f(x)$ неприводим над $\mathbb{Z}$. (То есть не представляется в виде произведения двух многочленов ненулевой степени с целыми коэффициентами)

\problem{5}
Докажите, что многочлен $x^{n-1}+x^{n-2}+\ldots+1$ неприводим тогда и только тогда, когда $n$ --- простое число.